\section{Hemalurgist's Weaponry}
\label{sec:hemalurgyWeapons}

While any metal weapon can be used to make a Hemalurgic Spike, a Hemalurgist may prefer to select tools designed for this purpose.

\subsection{Weapon Traits}

Some of the weapons presented in this chapter have new traits. You can find their definitions listed below: \\

\scriptedOption{Brutal}{
	Attacks made with this weapon gain an advantage to their damage dice.
}


\vspace{1cm}
\subsection{Weapons}


\begin{tabular}{l l l l l l}
	Skill 
	&Type
	& Damage
	& Range
	& Traits
	& Expert Traits 
	\\ \hline 
	
	Light Weaponry
	& Needle
	& 0-1 keen
	& Melee
	& Discreet, Subtle
	& Unique: Precise
	\\
	
	Light Weaponry
	& Small Spike
	& 1d4 keen
	& Melee
	& Quickdraw
	& Offhand
	\\
	
	Light Weaponry
	& Needle Blower
	& 0-1 keen
	& Ranged [30/120]
	& Quickdraw
	& Quiet
	\\
	
	Heavy Weaponry
	& Big Spike
	& 1d6 keen
	& Melee
	& Pierce
	& Offhand
	\\
	
	Heavy Weaponry
	& Hammer and Nail
	& 1d8 keen
	& Melee
	& Two-Handed, Deadly
	& Unique: Brutal
	\\
	
	Heavy Weaponry
	& Spike Thrower
	& 1d6 keen
	& Ranged [150/600]
	& Loaded[1], Two-Handed
	& Restraining
	\\
	
	
	
\end{tabular}



\newpage

\scriptedOption{Precise}{
	Attacks made with this weapon gain an advantage to determine if the attack hit the target.
}

\scriptedOption{Quiet}{
	Attacks made with this weapon do not alert the enemies about your position.
}

\scriptedOption{Restraining}{
	Spend \iconOpp{} from the attack with this weapon to make the target you hit Restrained.
}

\scriptedOption{Subtle}{
	While attacking with this weapon you may choose not to add your skill bonus to the damage dealt.
}